\introduction % Структурный элемент: ВВЕДЕНИЕ

В современных условиях программные системы всё чаще функционируют в
распределённой среде, где корректность их работы определяется не только
правильностью реализации отдельных модулей, но и согласованностью общего
поведения системы с формальной моделью. Для критичных и ответственных
программных решений особое значение приобретает применение методов,
позволяющих подтверждать соответствие исходного кода формальной
спецификации, выявлять расхождения между моделью и реализацией, а также
локализовывать ошибки на ранних этапах жизненного цикла программного
обеспечения.

В рамках данной выпускной квалификационной работы рассматривается методика
верификации соответствия TLA+ спецификации и исходного кода, основанная на
сопоставлении наблюдаемого исполнения распределённой реализации с её
формальной спецификацией. Предлагаемый подход использует
инструментирование программного кода, построение локальных и глобальных
трасс исполнения, а также последующую проверку согласованности полученной
трассы со спецификацией средствами модельной проверки. Такой подход
позволяет перейти от анализа низкоуровневых операций реализации к анализу
её поведения в терминах состояний и переходов, заданных в спецификации.

Организационно-правовой аспект данной работы связан прежде всего с тем,
что разработка и применение подобных методик затрагивают сферу
информационных технологий, защиты информации и безопасной разработки
программного обеспечения. При формировании трасс исполнения, обработке
событий, агрегировании журналов и анализе результатов верификации
необходимо учитывать требования законодательства Российской Федерации и
нормативных документов, регулирующих обращение с информацией, обеспечение
её целостности и конфиденциальности, а также требования к процессам
создания безопасного программного обеспечения.

Следовательно, организационно-правовая часть работы направлена на анализ
нормативных правовых актов и нормативных документов Российской Федерации,
которые задают общие принципы обращения с информацией, определяют
требования к её защите и формируют правовые и организационные основания
для применения методов верификации и безопасной разработки программных
средств. Рассмотрение указанных документов позволяет показать, что
разрабатываемая методика не противоречит действующему законодательству
Российской Федерации и соответствует современным подходам к обеспечению
надёжности и безопасности программных систем.

Основные нормативные правовые акты и нормативные документы,
рассматриваемые в данной организационно-правовой части, следующие:
\begin{itemize}
    \item Конституция Российской Федерации;
    \item Федеральный закон от 27.07.2006 №149-ФЗ <<Об информации, информационных технологиях и о защите информации>>;
    \item Указ Президента Российской Федерации от 05.12.2016 №646 <<Об утверждении Доктрины информационной безопасности Российской Федерации>>;
    \item ГОСТ Р 56939--2024 <<Защита информации. Разработка безопасного программного обеспечения. Общие требования>>;
    \item Приказ ФСТЭК России от 01.12.2023 №240 <<Об утверждении Порядка проведения сертификации процессов безопасной разработки программного обеспечения средств защиты информации>>.
\end{itemize}

Анализ указанных нормативных правовых актов и нормативных документов
позволяет определить правовой контекст разработки методики верификации
соответствия TLA+ спецификации и исходного кода, а также обосновать её
значение как инструмента повышения корректности, надёжности и
информационной безопасности программных систем.

